%!TEX root =  main.tex


\section{Frequently Asked Questions}
\label{asec:faq}


\para{Question 1.} Why does \alg{} not attain the maximum theoretical latency improvement of 50\% in the experiments?

\para{Answer.}
The theoretical upper bound of a 50\% latency reduction assumes an idealized setting in which drafting and verification times are perfectly matched throughout the execution of \alg{}. In practice, both stages exhibit stochastic variation due to system overheads, resource contention, and time multiplexing, making the precise temporal alignment required to attain this bound infeasible in real-world deployments.


\para{Question 2.} Why does \alg{} maintain two batches instead of more? Why is the ``Batch Parallelism size'' fixed to 2?

\para{Answer.} 
The design of \alg{} is intrinsically tied to the two-stage structure of speculative decoding (SD): drafting and verification. Compared to autoregressive decoding, the additional drafting stage introduces extra latency, which can become prohibitive when either the number of draft tokens per step is large or the draft model is computationally expensive. \alg{} mitigates this overhead by overlapping the drafting of one set of requests with the verification of another. This design naturally requires exactly two batches: one for drafting and one for verification. Introducing additional batches would not produce further parallelism benefits, as there are no additional stages to overlap, and would instead incur unnecessary scheduling and management overhead.


\para{Question 3.} Why does \alg{} not incorporate prefill requests in the parallelism?

\para{Answer.}
The current design of \alg{} assumes that chunked prefill is disabled. Under this configuration, vLLM completes the prefill stage for new requests in a dedicated step before transitioning them to decoding in subsequent steps, isolating prefill requests from the speculative decoding steps in which \alg{} operates. As support for chunked prefill is a planned extension, future versions of \alg{} will integrate prefill requests into its parallel execution workflow. \\